% =====================================================================
%  CARNET D'ENTRAINEMENT — FICHE 6 — THÈME 4
%  Niveau MOYEN — Corrigé
% =====================================================================
\documentclass[11pt,a4paper]{article}
\usepackage[utf8]{inputenc}
\usepackage[T1]{fontenc}
\usepackage{lmodern}
\usepackage[french]{babel}
\usepackage{amsmath,amssymb,mathtools}
\usepackage{icomma}
\usepackage[version=4]{mhchem}
\usepackage{siunitx}
\sisetup{output-decimal-marker={,}}
\usepackage{xcolor}
\usepackage{colortbl}
\usepackage{tikz}
\usepackage[most]{tcolorbox}
\usepackage{enumitem}
\usepackage{array}
\usepackage{booktabs}
\usepackage{fancyhdr}
\usepackage[hidelinks]{hyperref}
\usetikzlibrary{arrows.meta,positioning,calc,shapes.geometric}
\usepackage[a4paper,margin=2.1cm,headheight=26pt,headsep=8pt]{geometry}

\definecolor{primary}{RGB}{146,95,160}
\definecolor{primarydark}{RGB}{92,52,104}
\definecolor{corcol}{RGB}{0,114,178}
\definecolor{astucecol}{RGB}{198,93,9}
\definecolor{atomC}{RGB}{80,80,88}
\definecolor{atomH}{RGB}{235,235,235}
\definecolor{atomO}{RGB}{220,55,45}
\definecolor{atomN}{RGB}{48,80,200}
\definecolor{atomF}{RGB}{150,210,120}
\definecolor{atomS}{RGB}{225,190,40}
\definecolor{atomP}{RGB}{225,140,40}
\definecolor{lonepair}{RGB}{120,94,200}
\definecolor{bondcol}{RGB}{60,60,70}

\newcommand{\ball}[5]{\shade[ball color=#2] (#1) circle (#3);\node[font=\footnotesize\bfseries,text=#4] at (#1){#5};}
\newcommand{\pbond}[2]{\draw[bondcol,line width=1.6pt] (#1)--(#2);}
\newcommand{\wbond}[2]{\draw[bondcol,fill=bondcol,draw=none] let \p1=($(#2)-(#1)$), \n1={atan2(\y1,\x1)} in (#1) -- ($(#2)+(\n1+90:0.12)$) -- ($(#2)+(\n1-90:0.12)$) -- cycle;}
\newcommand{\hbond}[2]{\draw[bondcol,line width=1pt,dash pattern=on 1.4pt off 1.8pt] (#1)--(#2);}
\newcommand{\lobe}[3]{\begin{scope}[shift={(#1)},rotate=#2]\shade[ball color=lonepair!55,opacity=0.9] (#3,0) ellipse (0.34 and 0.20);\end{scope}}
\newcommand{\AX}[2]{\ensuremath{\mathrm{AX_{#1}E_{#2}}}}

\tcbset{boxbase/.style={enhanced,breakable,boxrule=0.9pt,arc=2.5pt,
   left=3mm,right=3mm,top=2.2mm,bottom=2.2mm,
   fonttitle=\bfseries,coltitle=white,toptitle=0.6mm,bottomtitle=0.6mm}}
\newtcolorbox[auto counter]{cor}[1][]{boxbase,colback=corcol!5,colframe=corcol,
   title={\textsc{Corrigé}~\thetcbcounter\ifx\relax#1\relax\relax\else\textmd{ —\itshape\ #1}\fi}}
\newtcolorbox{astucebox}[1][]{boxbase,colback=astucecol!8,colframe=astucecol,
   title={\textsc{Astuce}\ifx\relax#1\relax\relax\else\textmd{ : \itshape #1}\fi}}

\pagestyle{fancy}\fancyhf{}
\fancyhead[L]{\small\textcolor{primarydark}{\textbf{Fiche 6 · Thème 4} — VSEPR et géométrie moléculaire}}
\fancyhead[R]{\small\textcolor{primarydark}{\textsc{Moyen · Corrigé}}}
\fancyfoot[C]{\small\thepage}
\renewcommand{\headrulewidth}{0.8pt}
\renewcommand{\headrule}{\hbox to\headwidth{\color{primary}\leaders\hrule height \headrulewidth\hfill}}
\setlength{\parindent}{0pt}\setlength{\parskip}{3pt}

\newcommand{\rep}{\textbf{\textcolor{corcol}{Réponse : }}}

\begin{document}
\begin{center}
{\Large\bfseries\color{primarydark}Thème 4 — Corrigé du niveau Moyen}
\end{center}
\medskip

\begin{cor}[de la structure de Lewis aux deux géométries]
\begin{enumerate}[label=\alph*),leftmargin=2em,itemsep=3pt]
  \item \ce{SF2} : $n=2$, $m=2 \Rightarrow$ \AX{2}{2}. Répulsion \textbf{tétraédrique} (4 sites) ;
        géométrie moléculaire \textbf{coudée} (comme l'eau). \rep \AX{2}{2}, tétraédrique / coudée.
  \item \ce{PH3} : $n=3$, $m=1 \Rightarrow$ \AX{3}{1}. Répulsion \textbf{tétraédrique} ; géométrie
        moléculaire \textbf{pyramide trigonale} (comme \ce{NH3}). \rep \AX{3}{1}, tétraédrique /
        pyramide trigonale.
  \item \ce{BF3} : $n=3$, $m=0 \Rightarrow$ \AX{3}{0}. Le bore reste à 6 électrons (exception à
        l'octet). Répulsion \textbf{triangulaire} ; géométrie moléculaire \textbf{triangulaire
        plane}. \rep \AX{3}{0}, triangulaire / triangulaire plane.
\end{enumerate}
\begin{center}
\begin{tikzpicture}[scale=0.85]
  % SF2 coudée
  \coordinate (S) at (0,0);
  \coordinate (Fa) at (-1.25,-0.85);
  \coordinate (Fb) at ( 1.25,-0.85);
  \pbond{S}{Fa}\pbond{S}{Fb}
  \lobe{S}{68}{0.9}\lobe{S}{112}{0.9}
  \ball{Fa}{atomF}{0.32}{black}{F}
  \ball{Fb}{atomF}{0.32}{black}{F}
  \ball{S}{atomS}{0.42}{black}{S}
  \node[font=\scriptsize,align=center] at (0,-1.75){\ce{SF2} : \AX{2}{2} coudée};
  % BF3 triangulaire plane
  \begin{scope}[xshift=6cm]
    \coordinate (B) at (0,0);
    \coordinate (Fa) at (90:1.35);
    \coordinate (Fb) at (210:1.35);
    \coordinate (Fc) at (330:1.35);
    \pbond{B}{Fa}\pbond{B}{Fb}\pbond{B}{Fc}
    \ball{Fa}{atomF}{0.32}{black}{F}
    \ball{Fb}{atomF}{0.32}{black}{F}
    \ball{Fc}{atomF}{0.32}{black}{F}
    \ball{B}{atomN!30!white}{0.4}{black}{B}
    \node[primarydark,font=\scriptsize] at (150:0.9){$120^\circ$};
    \node[font=\scriptsize,align=center] at (0,-1.75){\ce{BF3} : \AX{3}{0} plane};
  \end{scope}
\end{tikzpicture}
\end{center}
\end{cor}

\begin{cor}[le cas du dioxyde de carbone \ce{CO2}]
\begin{enumerate}[label=\alph*),leftmargin=2em,itemsep=1pt]
  \item \rep Le carbone ne porte \textbf{aucun} doublet libre ($m=0$) : ses 4 électrons de valence
        engagés servent tous aux deux doubles liaisons.
  \item Chaque double liaison \ce{C=O} ne compte que pour \textbf{un} site de répulsion. Le carbone
        a donc $n=2$ atomes liés et $m=0$ : \rep type \AX{2}{0} (2 sites, pas 4).
  \item \rep Géométrie \textbf{linéaire}, angle \ce{O-C-O} $= \mathbf{180^\circ}$.
\end{enumerate}
\begin{center}
\begin{tikzpicture}[scale=0.9]
  \coordinate (C) at (0,0);
  \coordinate (O1) at (-1.9,0);
  \coordinate (O2) at ( 1.9,0);
  \draw[bondcol,line width=1.4pt] ($(C)+(-0.5,0.08)$)--($(O1)+(0.5,0.08)$);
  \draw[bondcol,line width=1.4pt] ($(C)+(-0.5,-0.08)$)--($(O1)+(0.5,-0.08)$);
  \draw[bondcol,line width=1.4pt] ($(C)+(0.5,0.08)$)--($(O2)+(-0.5,0.08)$);
  \draw[bondcol,line width=1.4pt] ($(C)+(0.5,-0.08)$)--($(O2)+(-0.5,-0.08)$);
  \ball{O1}{atomO}{0.45}{white}{O}
  \ball{O2}{atomO}{0.45}{white}{O}
  \ball{C}{atomC}{0.45}{white}{C}
  \draw[primarydark,thick,<->,>=Stealth] (-0.9,-0.8)--(0.9,-0.8);
  \node[primarydark,font=\scriptsize,fill=white,inner sep=1pt] at (0,-0.8){$180^\circ$};
\end{tikzpicture}
\end{center}
\end{cor}

\begin{cor}[la fermeture des angles]
\begin{enumerate}[label=\alph*),leftmargin=2em,itemsep=1pt]
  \item \rep \ce{CH4} ($109{,}5^\circ$) $>$ \ce{NH3} ($\approx 107^\circ$) $>$ \ce{H2O}
        ($\approx 104{,}5^\circ$).
  \item \rep Les doublets libres repoussent \emph{plus fort} que les doublets liants : plus l'atome
        central en porte, plus il comprime les liaisons et referme l'angle.
\end{enumerate}
\end{cor}

\begin{cor}[remplir un tableau de géométries]
\begin{center}
\renewcommand{\arraystretch}{1.25}
\footnotesize
\begin{tabular}{l c l}
\toprule
\textbf{Espèce} & \textbf{Type} & \textbf{Géométrie moléculaire}\\
\midrule
\ce{NH3} & \AX{3}{1} & pyramide trigonale\\
\ce{SO4^2-} & \AX{4}{0} & tétraédrique\\
\ce{N2} & \AX{1}{0}$^\star$ & linéaire (diatomique)\\
\ce{XeCl2} & \AX{2}{3} & linéaire\\
\ce{SF6} & \AX{6}{0} & octaédrique\\
\ce{BeCl2} & \AX{2}{0} & linéaire\\
\ce{ICl2^-} & \AX{2}{3} & linéaire\\
\bottomrule
\end{tabular}
\end{center}
\footnotesize $^\star$\ce{N2} (\ce{N#N}) est diatomique : deux atomes alignés, la molécule est
nécessairement \textbf{linéaire}.\par
\normalsize
\rep \ce{NH3} pyramide trigonale ; \ce{SO4^2-} tétraédrique ; \ce{N2} linéaire ; \ce{XeCl2}
linéaire (\AX{2}{3}) ; \ce{SF6} octaédrique ; \ce{BeCl2} linéaire (\AX{2}{0}) ; \ce{ICl2^-}
linéaire (\AX{2}{3}). Remarque : \ce{XeCl2} et \ce{ICl2^-} ont trois doublets libres équatoriaux,
mais leurs deux liaisons restent \emph{axiales} donc à $180^\circ$ : la molécule est linéaire.
\end{cor}

\begin{cor}[octet étendu et géométrie]
\begin{enumerate}[label=\alph*),leftmargin=2em,itemsep=1pt]
  \item \ce{PCl5} : $n=5$, $m=0 \Rightarrow$ \rep \AX{5}{0}, \textbf{bipyramide trigonale}.
  \item \ce{SF6} : $n=6$, $m=0 \Rightarrow$ \rep \AX{6}{0}, \textbf{octaédrique}.
  \item \ce{XeF4} : $n=4$, $m=2 \Rightarrow$ \rep \AX{4}{2}, \textbf{plan carré} (les 2 doublets
        libres se placent en positions opposées, au-dessus et au-dessous du plan).
\end{enumerate}
\end{cor}

\begin{astucebox}[atome central de période $\geq 3$]
Seuls les éléments à partir de la 3\textsuperscript{e} période (\ce{P}, \ce{S}, \ce{Cl}, \ce{Br},
\ce{I}, \ce{Xe}\ldots) peuvent porter \emph{plus} de 4 sites de répulsion (5 ou 6) grâce à l'octet
étendu : bipyramide trigonale et octaèdre leur sont réservés.
\end{astucebox}

\end{document}
